\frontmatter% Use roman page numbering style (i, ii, iii, iv...) for the pre-cont\begin{spacing}{1.5}
\addchaptertocentry{CANDIDATE DECLARATION}
\vspace{10 pt}
I hereby declare that work presented in the report entitled "\textbf{When Fine-Tuning Fails: Lessons from MS MARCO Passage Ranking}", submitted towards the fulfillment of MASTER'S THESIS report of M.Tech at Indian Institute of Information Technology Allahabad, is an authenticated original work carried out under supervision of \textbf{Dr. Muneendra Ojha}. Due Acknowledgements have been made in the text to all other material used. the project was done in full compliance with the requirements and constraints of the prescribed curriculum. 
\end{spacing}

\begin{spacing}{2.0}
\begin{flushright}
    \begin{minipage}{0.5\textwidth}
        \flushright \vspace{60 pt}
        \underline{\hspace{6cm}} \\
        \makebox[6cm]{\textbf{Manu Pande - MML2023005}} \\[80pt]declare that work presented in the report entitled "\textbf{When Fine-Tuning Fails: Lessons from MS MARCO Passage Ranking}", submitted towards the fulfillment of MASTER'S THESIS report of M.Tech at Indian Institute of Information Technology Allahabad, is an authenticated original work carried out under supervision of \textbf{Dr. Muneendra Ojha}. Due Acknowledgements have been made in the text to all other material used. the project was done in full compliance with the requirements and constraints of the prescribed curriculum.pages

\pagestyle{plain} % Default to the plain heading style until the thesis style is called for the body content
%----------------------------------------------------------------------------------------
%	TITLE PAGE
%----------------------------------------------------------------------------------------

\begin{titlepage}
\begin{center}
\begin{spacing}{1.2}

{\huge \bfseries When Fine-Tuning Fails: Lessons from MS MARCO Passage Ranking\par}  % Thesis title

\vspace {5mm}
\textit{A thesis submitted in partial fulfillment of the requirements\\for the award of the degree of} 

\vspace{7mm}
\textsc{\huge MASTER OF TECHNOLOGY}

\vspace {3mm}
\begin{figure}[htp]
    \centering
    \includegraphics[scale=0.28]{./IIITA_Logo.png}
\end{figure}
% ************************************
\begin{minipage}[t]{0.5\textwidth}
    \begin{flushleft} \large
        \textit{By:-} \\%[2mm]
            \textsc{MANU PANDE} % Student Names
    \end{flushleft}
\end{minipage}
\begin{minipage}[t]{0.45\textwidth}
    \begin{flushright} \large
        \textit{Enrollment No.} \\
            \textsc{MML2023005} % Student Enrollment numbers
    \end{flushright}
\end{minipage}\\[1cm]
% *************************************

\textit{Under the Supervision of}\\[2mm]
\textsc{\Large Dr. Muneendra Ojha}\\% Supervisor's Name

%\textit{Under the Co-Supervision of}\\[2mm]
%\textsc{\Large Prof. Manish Goswami}\\% Supervisor's Name
\vspace{7mm}
\textit{to the}\\[2mm]
\textsc{\Large Department of Information Technology}\\ % Department

\vspace{8mm}
\begin{hindi}
    \textsc{\Large {भारतीय सूचना प्रौद्योगिकी संस्थान, इलाहाबाद}} \\
\end{hindi}
\vspace{3mm}
\textsc{\Large Indian Institute of Information Technology, Allahabad} % university

\vspace{5mm}
{\fontsize{14}{14}\selectfont \text{June , 2025}}


\end{spacing}
\end{center}
\end{titlepage}



%----------------------------------------------------------------------------------------
%	DECLARATION
%----------------------------------------------------------------------------------------
\checktoopen
\begin{figure}[htp]
    % \centering
    \includegraphics[height=3cm,keepaspectratio]{./IIITA_Header.png}
\end{figure}
% \string\setblankpagestyle \space
\thispagestyle{empty}
\vspace{1mm}

\begin{center}
    {\large\bfseries CANDIDATE DECLARATION}
\end{center}

\begin{spacing}{1.5}
\addchaptertocentry{CANDIDATE DECLARATION}
\vspace{10 pt}
I hereby declare that work presented in the report entitled “\textbf{Design and Optimization of GO1\_GO2 Level Shifter by Analysing Across Multiple PVT Conditions}”, submitted towards the fulfillment of MASTER’S THESIS report of M.Tech at Indian Institute of Information Technology Allahabad, is an authenticated original work carried out under supervision of \textbf{Dr. Boddepalli Santhi Bhushan}. Due Acknowledgements have been made in the text to all other material used. the project was done in full compliance with the requirements and constraints of the prescribed curriculum. 
\end{spacing}

\begin{spacing}{2.0}
\begin{flushright}
    \begin{minipage}{0.5\textwidth}
        \flushright \vspace{60 pt}
        \underline{\hspace{6cm}} \\
        \makebox[6cm]{\textbf{Koneti Revanth Kumar - MMI2023011}} \\[80pt]
    \end{minipage}
\end{flushright}
\end{spacing}
\newpage

%----------------------------------------------------------------------------------------
%	CERTIFICATE
%----------------------------------------------------------------------------------------
\checktoopen
\begin{figure}[htp]
    % \centering
    \includegraphics[height=3cm,keepaspectratio]{./IIITA_Header.png}
\end{figure}
% \string\setblankpagestyle \space
\thispagestyle{empty}
\vspace*{.06\textheight}

\begin{spacing}{1.5}
\addchaptertocentry{CERTIFICATE FROM SUPERVISORS}
\begin{center}
    {\centering\large\bfseries CERTIFICATE FROM SUPERVISORS\par\vspace{10pt}}
\end{center}

\noindent It is certified that the work contained in the thesis titled \enquote{\textbf{When Fine-Tuning Fails: Lessons from MS MARCO Passage Ranking}} by \textbf{Manu Pande} has been carried out under supervision of \textbf{Dr. Muneendra Ojha} and that this work has not been submitted elsewhere for a degree.

\vspace{3.5cm}

\hfill\begin{minipage}{7.5cm}
    \begin{spacing}{1.2}
        \par
        \rule{\textwidth}{0.2pt}\\
        {Dr. Muneendra Ojha} \par
        {\deptname}  \par
        IIIT Allahabad \par
    \end{spacing}
\end{minipage}


\end{spacing}
% \end{certificate}
\cleardoublepage

\begin{figure}[htp]
    % \centering
    \includegraphics[height=3cm,keepaspectratio]{./IIITA_Header.png}
\end{figure}
% \string\setblankpagestyle \space
\thispagestyle{empty}
\vspace*{.06\textheight}

\begin{spacing}{1.5}
\addchaptertocentry{CERIFICATE OF APPROVAL}
\begin{center}
    {\centering\large\bfseries CERTIFICATE OF APPROVAL \par\vspace{10pt}}
\end{center}

 This thesis entitled \textbf{When Fine-Tuning Fails: Lessons from MS MARCO Passage Ranking}  by \textbf{Manu Pande} (MML2023005) is approved for the degree of Master's thesis at IIIT Allahabad  It is understood that by this approval, the undersigned does not necessarily endorse or approve any statement made, opinion expressed, or conclusion drawn
therein but approves the thesis only for the purpose for which it is submitted."\\[40 pt]

 Signature and name of the committee members (on final examination and approval of the thesis): \\[10 pt]
 \begin{flushleft}
     \begin{enumerate}
         \item Dr. Boddepalli Santhi Bhushan \\[25pt]
         \item Dr. Kavindra Kandpal\\[25 pt]
         \item Dr. Anand Kumar Tiwari \\[25 pt]

     \end{enumerate}
 \end{flushleft}
 \vspace{20 pt}
 \begin{flushright}
     \begin{minipage}{0.5\textwidth}
        \flushright \vspace{60 pt}
        \underline{\hspace{6cm}} \\
        \makebox[6cm]{\textbf{Dean(A\&R)}} \\[80pt]
    \end{minipage}
 \end{flushright}
\end{spacing} 

\newpage
\addchaptertocentry{PLAGIARISM REPORT}
\includepdf[pages=-]{plagrevanth.pdf}
\newpage
\begin{figure}[htp]
    % \centering
    \includegraphics[height=3cm,keepaspectratio]{./IIITA_Header.png}
\end{figure}
% \string\setblankpagestyle \space
\thispagestyle{empty}
\vspace{1mm}

\begin{center}
    {\large\bfseries ACKNOWLEDGEMENT}
\end{center}

\begin{spacing}{1.5}
\addchaptertocentry{ACKNOWLEDGEMENT}
I am thankful to my project supervisor, \textbf{Dr. Muneendra Ojha} for the guidance, support, and invaluable feedback throughout the research process. Their expertise, encouragement, and patience have been instrumental in the completion of this thesis.\\[10pt] 
 I am grateful to Dr. Sanjay Singh Sir, Head of the Department of Electronics and Communication Engineering, and also my panel members Dr Kavindra Kandpal and Prof Anand Kumar for their guidance and support.\\[10pt]
 I would also like to thank the Indian Institute of Information Technology, Allahabad for providing the necessary resources and facilities to conduct this research. I am deeply grateful to the participants who generously shared their time and personal information to make this study possible. \\[30 pt]
\end{spacing}

\begin{spacing}{2.0}
\begin{flushright}
    \begin{minipage}{0.5\textwidth}
        \flushright \vspace{60 pt}
        \underline{\hspace{6cm}} \\
        \makebox[6cm]{\textbf{Koneti Revanth Kumar - MMI2023011}} \\[80pt]
    \end{minipage}
\end{flushright}
\end{spacing}
\newpage


\begin{center}
    {\large\bfseries ABSTRACT}
\end{center}
\begin{spacing}{1.5}
\addchaptertocentry{ABSTRACT} % Add the abstract to the table of contents
     This thesis presents the design, characterization, and optimization of a dual-gate (GO1–GO2) level shifter for multi-power domain System-on-Chip (SoC) applications. Power consumption minimization is crucial in modern semiconductor design due to increasing integration density and portable application demands. Level shifters are essential interface circuits that facilitate reliable signal communication between different voltage domains, ensuring logic integrity while crossing domain boundaries. Conventional level shifters are divided into cross-coupled and Current Mirror Based level shifters. This research aims to enhance a cross-coupled topology through dual-gate implementation. The proposed level shifter was analyzed using Spectre simulation, deterministic worst-case analysis, statistical Monte Carlo simulations, and High Sigma Verification (HSV) methodology. Circuit optimization focused on minimizing propagation delay through strategic transistor sizing adjustments while maintaining signal integrity. Modifications to the pull-up network reduced drive strength, improving cell\_rise and cell\_fall timing characteristics without compromising contention resolution capabilities. These optimizations substantially reduced delays while minimally impacting leakage power consumption and silicon area requirements. The optimized design demonstrates enhanced performance across all operating conditions while maintaining robustness against process variability. Comparative analysis with existing solutions shows significant improvements in speed-power metrics, making the proposed level shifter suitable for high-performance, low-power SoC applications.
\end{spacing}


\newpage


%----------------------------------------------------------------------------------------
%	LIST OF CONTENTS/FIGURES/TABLES PAGES
%----------------------------------------------------------------------------------------
\addchaptertocentry{Table of Contents}
\tableofcontents % Prints the main table of contents
\listoffigures
\listoftables

%----------------------------------------------------------------------------------------
%	ABBREVIATIONS
%----------------------------------------------------------------------------------------
% \chapter*{List of Abbreviations}
% \addcontentsline{toc}{chapter}{List of Abbreviations}
% \begin{tabular}{@{}ll}
%     Machine Learning Models\\
%     $E_{m}$ & Energy level\\
%     $S_{m}$ & MMSE estimation \\
%     SNR\\
%     Ambient Back scattering Communication \\
%     SVM, RF, DT, XGB \\
%     BER and reflection Coefficient \\
%     Tag decode \\
%     % Add more as needed
% \end{tabular}

%----------------------------------------------------------------------------------------
%	PHYSICAL CONSTANTS/OTHER DEFINITIONS
%----------------------------------------------------------------------------------------

% \begin{constants}{lr@{${}={}$}l} % The list of physical constants is a three column table

% % The \SI{}{} command is provided by the siunitx package, see its documentation for instructions on how to use it

% Speed of Light & $c_{0}$ & \SI{2.99792458e8}{\meter\per\second} (exact)\\
% Constant Name & $Symbol$ & $Constant Value$ with units\\

% \end{constants}

%----------------------------------------------------------------------------------------
%	SYMBOLS
%----------------------------------------------------------------------------------------

% \begin{symbols}{lll} % Include a list of Symbols (a three column table)

% $a$ & distance & \si{\meter} \\
% $P$ & power & \si{\watt} (\si{\joule\per\second}) \\
% %Symbol & Name & Unit \\

% \addlinespace % Gap to separate the Roman symbols from the Greek

% $\omega$ & angular frequency & \si{\radian} \\

% \end{symbols}

%----------------------------------------------------------------------------------------
%	DEDICATION
%----------------------------------------------------------------------------------------